\documentclass[11pt,a4paper]{article}

% ==================================================
% Paquetes
% ==================================================
\usepackage[utf8]{inputenc}
\usepackage[T1]{fontenc}
\usepackage[english]{babel}
\usepackage{lmodern}
\usepackage[top=1.8cm,bottom=1.8cm,left=2cm,right=2cm]{geometry}
\usepackage{xcolor}
\usepackage{enumitem}
\usepackage{titlesec}
\usepackage{tabularx}
\usepackage{microtype}
\usepackage{hyperref}

\definecolor{azul}{RGB}{28,70,110}
\definecolor{gris}{RGB}{90,90,90}

\hypersetup{colorlinks=true, urlcolor=azul, linkcolor=azul,
  pdftitle={Henry Barrera - CV}, pdfauthor={Henry Barrera}}

\pagestyle{empty}
\setlength{\parindent}{0pt}
\setlist[itemize]{leftmargin=1.2em, itemsep=1pt, topsep=2pt}

\titleformat{\section}{\large\scshape\color{azul}}{}{0em}{}[\vspace{-0.6ex}\color{azul}\rule{\textwidth}{0.4pt}]
\titlespacing*{\section}{0pt}{1.6ex}{1ex}

% Telefono: solo aparece si existe cv/phone.tex (archivo local, ignorado por git).

% \entry{titulo}{fechas}{lugar / contexto}
\newcommand{\entry}[3]{%
  \begin{tabularx}{\textwidth}{@{}>{\raggedright\arraybackslash}X@{\hspace{1em}}r@{}}
    \textbf{#1} & {\color{gris}#2}
  \end{tabularx}\\
  {\color{gris}\small #3}\par\vspace{2pt}}

\begin{document}

% ==================================================
% Encabezado
% ==================================================
\begin{center}
  {\LARGE\scshape Henry Barrera}\\[4pt]
  {\color{gris} M.Sc. student in Applied Mathematics \textperiodcentered{} Bogot\'a, Colombia}\\[3pt]
  \small
  \href{mailto:hbarrera@unal.edu.co}{hbarrera@unal.edu.co}
  \IfFileExists{phone.tex}{\textperiodcentered{} \input{phone.tex}}{}
  \textperiodcentered{} \href{https://henry2899.github.io}{henry2899.github.io}
  \textperiodcentered{} \href{https://github.com/Henry2899}{GitHub}
  \textperiodcentered{} \href{https://www.linkedin.com/in/henry-barrera-3422981aa}{LinkedIn}
\end{center}

% ==================================================
\section{Profile}
% ==================================================
Mathematician from the Universidad Nacional de Colombia, currently pursuing an M.Sc.\ in Applied
Mathematics with a thesis on discrete optimization and quantum computing. I build mathematical models
of real problems, solve them numerically, estimate their parameters from data and assess how far the
results can be trusted. Experience with integer programming, differential equation models, Bayesian
inference and uncertainty quantification, mostly in Python and Julia.

% ==================================================
\section{Education}
% ==================================================
\entry{M.Sc.\ in Applied Mathematics}{2026 -- expected 2028}{Universidad Nacional de Colombia, Bogot\'a}
Thesis (in progress): discrete optimization and quantum computing, with the Traveling Salesman Problem
as reference problem.

\vspace{4pt}
\entry{Matem\'atico (professional degree in Mathematics)}{2019 -- 2026}{Universidad Nacional de Colombia, Bogot\'a}

% ==================================================
\section{Research experience}
% ==================================================
\entry{Student research assistant -- HERMES 500 project}{2026}{Universidad Nacional de Colombia
  \textperiodcentered{} Supervisor: Prof.\ Vladimir Angulo Castillo \textperiodcentered{}
  \href{https://github.com/Henry2899/ic50-uncertainty-toolkit}{toolkit (synthetic data)}}
\begin{itemize}
  \item Built a reproducible Python workflow for the DPPH antioxidant assay; estimated IC$_{50}$
        values by ordinary and weighted least squares.
  \item Quantified the uncertainty of IC$_{50}$ with three independent methods: delta method,
        Monte Carlo simulation and nonparametric bootstrap.
  \item Designed a Mamdani fuzzy inference system (2 inputs, 9 rules) that combines antioxidant
        potency and experimental precision into a 0--100 ranking index.
  \item Audited a fuzzy model proposed within the project and reported 8 methodological issues
        (interpolation and extrapolation assumptions, membership functions, incomplete rules).
  \item Analysed total phenolic content across five plant species with ANOVA, Tukey comparisons
        and assumption checks (Shapiro--Wilk, Levene).
\end{itemize}

% ==================================================
\section{Projects}
% ==================================================
\entry{Discrete optimization and quantum computing (M.Sc.\ thesis)}{2026 -- present}{Based on Bochkarev
  et al., \emph{European Journal of Operational Research} (2025)}
\begin{itemize}
  \item Modelled the TSP with the Dantzig--Fulkerson--Johnson, Miller--Tucker--Zemlin and
        quadratic-assignment formulations and compared their size.
  \item Solved benchmark instances with Gurobi; next stage: QUBO reformulation for quantum
        annealing and variational quantum algorithms.
\end{itemize}

\entry{Siting distributed solar generation on an urban power grid}{2026}{Exploratory thesis project,
  paused \textperiodcentered{} Python, NetworkX \textperiodcentered{} \href{https://github.com/Henry2899/pvgraph}{repository}}
\begin{itemize}
  \item Built a Python package (44 automated tests) for the distribution grid of Yopal (Casanare):
        graph construction, data validation, and voltage (LinDistFlow), loading, loss and radiality
        checks on a reproducible 322-node synthetic network.
  \item Formulated a mixed-integer model for siting photovoltaic generation and a demand-forecasting
        stage; not implemented before the project was paused.
\end{itemize}

\entry{Simulated annealing for the Traveling Salesman Problem}{2026}{Markov Chains course
  \textperiodcentered{} with L.\ Torres \textperiodcentered{}
  \href{https://github.com/Henry2899/cadenas-markov-aplicaciones}{repository}}
\begin{itemize}
  \item Simulated annealing studied as a Markov chain on the space of tours.
\end{itemize}

\entry{An HSIR model of dengue in Neiva, 2024}{2025}{Numerical Analysis course (Prof.\ Juan Carlos
  Galvis) \textperiodcentered{} team of three \textperiodcentered{} Python, Julia \textperiodcentered{}
  \href{https://github.com/Henry2899/dengue-hsir-neiva}{repository}}
\begin{itemize}
  \item Built the dataset from the weekly epidemiological bulletins of Huila: wrote extraction
        scripts for a format that changed during the year (PyMuPDF, pdfplumber, regular expressions,
        OCR with pytesseract) and completed the data by manual transcription where the automatic
        extraction was unreliable.
  \item The team fitted the model by least squares and Bayesian inference (Turing.jl) and compared
        five numerical integrators.
\end{itemize}

\entry{Orbital modelling of the Starlink-1106 satellite}{2025}{Mathematical Modelling I course
  \textperiodcentered{} team of four \textperiodcentered{} Julia \textperiodcentered{}
  \href{https://github.com/Henry2899/starlink-1106-orbit}{repository}}
\begin{itemize}
  \item Four orbital models of increasing complexity (Earth oblateness, atmospheric drag, lunar and
        solar perturbations, relativistic corrections) fitted to real ephemeris data with
        Nelder--Mead and Metropolis--Hastings.
  \item The best model reduced the least-squares residual from $1.96\times10^{9}$ to
        $2.87\times10^{6}$.
\end{itemize}

% ==================================================
\section{Talks and seminars}
% ==================================================
\entry{Experimental mathematics in combinatorics: cases, symmetries and computation}{Dec 2026}{Combinata:
  Encuentro de Combinatoria, \'Algebra, Teor\'ia de N\'umeros y sus Aplicaciones \textperiodcentered{}
  UPTC, Duitama \textperiodcentered{} in Spanish}
\entry{Chasing the Shortest Route: Inside the Traveling Salesman Problem}{Nov 2026}{Seminario de
  Matem\'atica Aplicada \textperiodcentered{} Universidad Nacional de Colombia}
\entry{The Lorenz attractor and the geometric Lorenz attractor}{Jul 2026}{Hyperbolic Dynamics
  seminar \textperiodcentered{} Universidad Nacional de Colombia}

% ==================================================
\section{Tutoring}
% ==================================================
Mathematics tutor, Universidad Nacional de Colombia (Feb -- Jun 2023) \textperiodcentered{}
Private tutor in mathematics, Bogot\'a (2023 -- present)

% ==================================================
\section{Skills}
% ==================================================
\begin{tabularx}{\textwidth}{@{}p{3.2cm}X@{}}
  Programming & Python (NumPy, SciPy, pandas, statsmodels, NetworkX, matplotlib, pytest),
                Julia (DifferentialEquations.jl, Turing.jl, Optim.jl), SQL \\
  Optimization & Integer programming with Gurobi (gurobipy), heuristics and metaheuristics \\
  Statistics & Regression (OLS/WLS), ANOVA, Bayesian inference and MCMC, Monte Carlo, bootstrap \\
  Other & \LaTeX, Git, Power BI, Excel \\
  Languages & Spanish (native), English (B1), Italian (beginner)
\end{tabularx}

% ==================================================
\section{Additional training}
% ==================================================
Python for Data Science (Pontificia Universidad Cat\'olica de Chile) \textperiodcentered{}
Data Analysis and Visualization (University at Buffalo) \textperiodcentered{}
Hypothesis Testing with Python and Excel (Tufts University) \textperiodcentered{}
Corporate Finance specialization (UNAM)

\end{document}
